\documentclass[UTF8,fontset=none,11pt,a4paper]{ctexart}
\usepackage{ipara-handbook}
\handbooksetup{DS-10 Hash 与直接定位}{408 · 数据结构 DS-10}{Key · Collision · Load · Expected Cost}
\begin{document}
\handbooktitle{Hash 与直接定位}{用槽位和冲突管理换取接近直接访问}{408 · 数据结构 Topic DS-10}

\begin{corebox}{母问题：映射不是定位，冲突怎样成为状态的一部分？}
散列的生成链是
\[
\text{Key}\to\text{Hash Mapping}\to\text{Collision}\to\text{Resolution}\to\text{Load State}\to\text{Expected Cost}.
\]
哈希函数只给出候选槽位，不保证唯一；开放定址沿探测序列寻找空位，链地址把同槽 key 放入桶链。查找失败必须按同一探测规则走到“从未使用”或完整探测，删除不能简单清空开放定址槽位，否则会截断后续 key 的探测链。
\end{corebox}
\begin{methodbox}{本册定位与 Owner 边界}
\begin{itemize}
\item \kw{Owns：}散列函数、装填因子、开放定址/链地址、删除标记、重散列与成功/失败查找语义。
\item \kw{Uses：}Atlas Foundation 与 DS-B02 workload 比较；精确等值查询优先，不承诺范围顺序。
\item \kw{Stops：}有序索引归 DS09；具体工程哈希表的并发、缓存和安全策略只作 Extension。
\end{itemize}
\end{methodbox}
\tableofcontents\newpage

\section{状态与冲突策略}
开放定址槽位至少有 Empty、Occupied、Deleted 三态；链地址槽位可用容器表示。装填因子 $\alpha=n/m$ 越高，冲突和探测长度通常越大；超过阈值应扩容并重新插入，不能直接把旧下标复制到新表。删除标记仍占探测路径，重散列时可被清除。
\begin{examplebox}{为什么删除后不能变 Empty}
若 $h(a)=h(b)$，先插入 $a$ 再插入 $b$，$b$ 可能探测到下一个槽。删除 $a$ 后若把槽清空，查找 $b$ 看到第一个 Empty 就会错误失败；Deleted 表示“继续探测，但此处可复用”。
\end{examplebox}

\subsection{哈希函数先看取值域与分布}
直接定址 $H(k)=k$ 只适合 key 值域小且连续的场景；常见除留余数法写成 $H(k)=k\bmod m$，表长 $m$ 通常选质数以减少周期性碰撞。折叠、平方取中等方法适合 key 的位分布不同的场景，但任何函数都只能给出候选地址，不能证明无冲突。分析失败查找的 ASL 时，入口数由 $H$ 的实际值域决定，而不是看到容量 $m$ 就直接除以 $m$。

\subsection{链地址与开放定址的两种成本}
链地址把同一槽位的记录挂在桶链上，装填因子可以超过 $1$；查找先计算桶，再在链内比较。均匀散列下，成功查找的平均链长约为 $1+\alpha/2$，失败查找约为 $\alpha$（具体常数随“插入前/插入后计数”口径变化，题目应按给定定义计算）。删除可直接摘除链节点，但要处理头节点和重复 key。

开放定址要求 $\alpha<1$，所有记录都在表内，冲突后按探测序列
\[
H_i(k)=\bigl(H(k)+i\bigr)\bmod m
\]
或二次探测、双重散列继续找槽。线性探测实现简单但会产生主聚集；二次探测减轻主聚集，却要检查表长与步长是否覆盖所有槽；双重散列的第二步长必须与 $m$ 互素，否则可能只访问部分槽位。探测序列一旦绕回起点仍未命中，就应报告满表/失败而不是继续循环。

\subsection{装填因子、删除与重散列}
插入时应区分“第一次 Deleted 槽位”和“真正 Empty 槽位”：可以记住第一个 Deleted 位置，但仍需继续探测，确认 key 不在后再复用它。扩容后每个 key 的 $H(k)\bmod m'$ 都可能改变，必须按新表长重新插入，不能把旧数组下标整体复制。缩容、长期 tombstone 堆积和链地址桶过长都是重散列触发条件；扩容阈值是工程策略，题目若未给定不能擅自当成固定常数。

\begin{methodbox}{哈希题的失败证明}
成功查找：按同一探测序列遇到目标。失败查找：开放定址遇到从未使用的 Empty，或已经完整探测一周；链地址扫描完桶链。Deleted 只能表示“此处曾经有记录”，不能作为开放定址的失败终点。把三种状态和终止条件写在草稿上，才能避免删除后误报不存在或满表死循环。
\end{methodbox}

\section{成本与边界}
均匀散列、合理装填因子和独立探测假设下，开放定址成功/失败查找期望接近常数；最坏可退化到 $O(m)$。哈希不保序，因此范围查询、前驱后继和有序遍历不能由平均 $O(1)$ 推出。

\section{Hash 作为辅助索引：把全局扫描改成局部证据}
Hash 的算法价值常常不是“单独存一张表”，而是为扫描过程维护已见集合、频次或对象到位置的反向索引。正确性仍来自外层算法不变量，期望复杂度则额外依赖哈希分布与装填控制。

\subsection{配对、分组与连续段}
两数之和的一遍扫描维护“此前已经出现的值 $\to$ 位置”。处理 $x$ 时先查询目标补数 $t-x$，命中后得到两个不同位置；再插入当前 $x$，从而不会在只出现一次时重复使用同一元素。若需要所有配对，重复值、输出去重和位置组合会改变接口，不能沿用“找到即返回”的状态。

异位词分组需要把等价对象压成同一规范 key：把每个字符串排序后作为 key，每项成本 $O(L\log L)$；固定小字符集时用频次数组作 key，可降为 $O(L+|\Sigma|)$，但字符集和整数编码必须固定。Hash 只负责把相同 key 聚到一起，“规范化是否保持且仅保持目标等价关系”才是分组正确性的核心。

最长连续整数段先把所有不同值放入集合，只从不存在 $x-1$ 的值 $x$ 开始向右扩展。这样每个值至多被一条真正的段扫描访问一次，均匀散列假设下期望 $O(n)$；若从每个值都向两侧扩展，会对长段反复扫描而退化。重复输入应在集合化阶段消除，极值处计算 $x-1$ 或 $x+1$ 还需处理整数溢出。

\subsection{对象到位置的反向索引}
当数组支持尾部 $O(1)$ 插入/删除时，Hash 可维护 \code{value -> index}：删除任意值时把数组末元素搬到空位，并同步更新其索引。这是 RandomizedSet、频率桶和缓存结构常用接口，但组合结构还必须维护数组、链表、频次桶等跨索引一致性；完整的 workload、失败原子性与 LRU/LFU 推导由 DS-I01 Own，本册只提供反向定位机制。

\begin{warnbox}{期望线性不是无条件线性}
上述应用把每次 Hash 操作按期望 $O(1)$ 计；对抗输入、退化哈希函数、重散列峰值或必须给出最坏界时，应改用有序树、排序或受控哈希策略，并重新报告成本向量。
\end{warnbox}

\section{代码与测试证据}
完整实现位于 \codepath{code/ds10_hash_table.hpp}，测试位于 \codepath{code/ds10_hash_table_test.cpp}。
\lstinputlisting[language=C++,firstline=12,lastline=108,firstnumber=1,numbers=left,caption={线性探测、删除标记与重散列}]{30_408/10_数据结构/10_Hash与直接定位/code/ds10_hash_table.hpp}
\lstinputlisting[language=C++,firstline=8,lastline=47,firstnumber=1,numbers=left,caption={冲突、删除后查找与满表边界测试}]{30_408/10_数据结构/10_Hash与直接定位/code/ds10_hash_table_test.cpp}
\begin{lstlisting}[language=bash]
clang++ -std=c++17 -Wall -Wextra -Werror -pedantic code/ds10_hash_table_test.cpp -o /tmp/ds10_test
/tmp/ds10_test
\end{lstlisting}

\section{做题控制协议}
先问查询是否精确等值，再写 $h(k)$、冲突处理、装填因子和失败终止条件。开放定址题画探测序列，删除保留 tombstone；扩容题重新计算每个 key 的槽位。若题目要求范围或排序，转交 DS09/DS11。
\begin{corebox}{压缩复原}
五问：key 映射到哪里？冲突后走哪条路径？槽位有几种状态？失败何时停止？装填变大如何修复？
\end{corebox}
\end{document}
